Constraint satisfaction credits its performance to two mechanisms that share the family's name:
\emph{ordering} --- choosing which variable to instantiate next (minimum remaining values, degree
tie-break) and which value to try first (least constraining value) --- and \emph{propagation} ---
pruning domains after each assignment (forward checking). Both are taught as the difference between
an exponential blow-up and a tractable search.

We ran five solvers --- chronological backtracking, {\tt bt\_mrv}, {\tt bt\_lcv},
{\tt bt\_fc\_mrv\_deg} (forward checking + MRV + degree), and min-conflicts local search
\citep{minton1992minimizing} --- over 10 instance sizes $\times$ 8 seeds at a 5\,s budget per
instance: 400 runs. Every run records whether it exhausted the budget, which turns out to matter
more than anything else in this section.

\subsection{The orderings do not change what gets solved}

Table~\ref{tab:csp-solved} gives the outcome that the ordering literature would most expect to
move, and it does not move. All four systematic solvers solve the same number of instances to
within one, and pairwise they disagree on the outcome of \textbf{1 of 80 instances} --- at $n = 25$,
seed 303, where {\tt bt\_lcv} fails and the other three succeed. That single disagreement runs
\emph{against} the heuristic. Adding MRV, LCV, or forward checking to chronological backtracking
changed which instances were solvable, at this budget, essentially not at all.

\begin{table}[t]
\centering
\small
\setlength{\tabcolsep}{3.5pt}
\caption{Left: outcomes over all 400 runs. Right: the 24 instances every solver solved, where node
counts are uncensored and therefore meaningful; medians there, means on the left. Paired Wilcoxon
against chronological backtracking, on node counts.}
\label{tab:csp-solved}
\begin{tabular}{@{}lrrrrrl@{}}
\toprule
& \multicolumn{3}{c}{\textbf{all 400 runs}} & \multicolumn{3}{c}{\textbf{24 commonly solved instances}} \\
\cmidrule(lr){2-4}\cmidrule(lr){5-7}
\textbf{Solver} & \textbf{solved} & \textbf{nodes} & \textbf{time} &
\textbf{nodes} & \textbf{time} & \textbf{vs.\ basic} \\
\midrule
Chronological backtracking & 27/80 & 444,111 & 3.15\,s & 13.5 & 0.09\,ms & --- \\
\quad + MRV / degree       & 27/80 & 322,997 & 3.02\,s & 13.5 & 0.14\,ms & $+0.0$, $p = 0.27$ \\
\quad + LCV                & 26/80 & 223,080 & 3.23\,s & 10.5 & 1.09\,ms & $-3.0$, $p = 10^{-4}$ \\
\quad + FC + MRV + deg.\ & 27/80 & 126,481 & 3.17\,s & \phantom{0}9.0 & 1.27\,ms & $-5.0$, $p < 10^{-4}$ \\
\midrule
Min-conflicts (repair) & 25/80 & \textbf{119} & \textbf{0.044\,s} & \textbf{1.0} & \textbf{0.08\,ms} & $-13.5$, $p < 10^{-4}$ \\
\bottomrule
\end{tabular}
\end{table}

\subsection{The orderings do reduce nodes --- and make the solver slower}

On the 24 instances every solver solved --- the only regime where a node count means what it is
supposed to mean --- the ordering heuristics behave exactly as folklore predicts, and the
improvement is statistically solid. Least-constraining-value removes a median of 3.0 nodes
($p = 10^{-4}$) and forward checking with MRV removes 5.0 ($p < 10^{-4}$), against a median of 13.5
for chronological backtracking. MRV on its own is the exception: a median difference of exactly
zero ($p = 0.27$).

Every one of them is slower. LCV takes $12.1\times$ and forward checking $14.1\times$ the wall clock
of chronological backtracking on the same instances (1.09\,ms and 1.27\,ms against 0.09\,ms), and
even MRV --- which saves no nodes at all --- costs $1.6\times$. The per-node price of computing the
ordering and running the propagation exceeds, by an order of magnitude, the value of the nodes it
saves. This is the same structure as \S\ref{sec:search}: the customary metric shows the eponymous
mechanism working, and it \emph{is} working, but the metric does not charge what it costs to run.

The one arm that improves both metrics simultaneously is the one that abandons systematic search.
Min-conflicts expands a median of 1.0 node at 0.08\,ms --- as fast as chronological backtracking
per attempt and $13.5$ nodes cheaper --- and over all 400 runs averages 119 nodes and 0.044\,s
against 444,111 nodes and 3.15\,s: $3{,}700\times$ fewer nodes and $72\times$ less time, for
25 solved instances against 27. Both solve rates are low in absolute terms --- $31\%$ and $34\%$ of
each solver's 80 runs, because most instances in this sweep are out of reach for every solver at a
5\,s budget --- but min-conflicts reaches $93\%$ of chronological backtracking's solve count for $0.03\%$ of its
search cost, by carrying one complete (initially infeasible) assignment forward and repairing it
rather than rebuilding a partial assignment prefix on every backtrack.

\subsection{Under censoring, the node metric inverts}

At $n \geq 30$ every systematic run exhausts the 5\,s budget and nothing is solved. Aggregate node
counts over such runs are widely reported anyway --- our own source artifact reports them --- and
they are worse than uninformative. When every solver stops at the same clock, its node count
measures \emph{how fast it grinds}, not how well it prunes:

\begin{center}
\small
\begin{tabular}{@{}lrrr@{}}
\toprule
\textbf{Solver} & \textbf{$n=30$} & \textbf{$n=40$} & \textbf{$n=50$} \\
\midrule
Chronological backtracking & 163,065 & 179,949 & 168,470 \\
\quad + MRV / degree       & \phantom{0}82,923 & 114,392 & 117,990 \\
\quad + LCV                & \phantom{0}64,761 & \phantom{0}60,027 & \phantom{0}35,103 \\
\quad + FC + MRV + degree  & \phantom{0}30,066 & \phantom{0}22,951 & \phantom{0}16,588 \\
\bottomrule
\end{tabular}
\\[2pt]
\footnotesize nodes per second on fully censored cells (all runs hit the budget; none solved)
\end{center}

Forward checking's apparent $3.5\times$ node advantage in the all-runs column of
Table~\ref{tab:csp-solved} is therefore not a pruning result at all: it is lower throughput
measured against a shared stopwatch, by $5.4\times$ at $n = 30$, $7.8\times$ at $n = 40$ and
$10.2\times$ at $n = 50$. Under censoring the sign of the metric's meaning
flips --- fewer nodes indicates a \emph{slower} solver, not a better-pruned tree --- and the two
cases are indistinguishable in the reported number. Any scaling curve of nodes against $n$ that
crosses into a censored regime without saying so is reporting throughput and calling it efficiency.

\subsection{Verdict}

Ordering and propagation do what they claim to do, measured in the unit the field uses, on the
instances where that unit is well defined. They also cost $12$--$14\times$ more wall clock than they
save at these sizes, change which instances are solvable in 1 case out of 80 (and that one
adversely), and produce node counts at larger sizes that measure the opposite of what they appear
to. The mechanism that actually changes the cost of solving these instances by orders of magnitude
is the one the family is not named for.
